% The supplement to the Linnik exponent paper (main.tex): the proofs of four standard results, and the details of the
% paper's Sections 7 and 8. Compile main.tex first: the references to the paper (\ref{M-...}) are read from main.aux.
\documentclass[11pt]{article}
\usepackage[T1]{fontenc}
% pdfLaTeX needs the input encoding and the glyph declarations below; XeTeX reads Unicode directly.
\ifdefined\XeTeXversion\else\usepackage[utf8]{inputenc}\fi
\usepackage{amsmath,amssymb,amsthm}
\usepackage{booktabs}
\usepackage{longtable}
\usepackage{xr-hyper}
\usepackage[hidelinks]{hyperref}
% References to the paper, compiled first: \ref{M-lem:prime-sum} is its Lemma~3.2.
\externaldocument[M-]{main}
% The Unicode symbols of the Lean statements quoted in Section S5.
\ifdefined\XeTeXversion\else
\DeclareUnicodeCharacter{2203}{\ensuremath{\exists}}
\DeclareUnicodeCharacter{2200}{\ensuremath{\forall}}
\DeclareUnicodeCharacter{211D}{\ensuremath{\mathbb R}}
\DeclareUnicodeCharacter{2115}{\ensuremath{\mathbb N}}
\DeclareUnicodeCharacter{2102}{\ensuremath{\mathbb C}}
\DeclareUnicodeCharacter{2227}{\ensuremath{\wedge}}
\DeclareUnicodeCharacter{2192}{\ensuremath{\to}}
\DeclareUnicodeCharacter{2261}{\ensuremath{\equiv}}
\DeclareUnicodeCharacter{2264}{\ensuremath{\le}}
\DeclareUnicodeCharacter{03C7}{\ensuremath{\chi}}
\fi


\newcommand{\lean}[1]{\texttt{#1}}
\newcommand{\file}[1]{\nolinkurl{#1}}
\emergencystretch=3em
\numberwithin{equation}{section}
\renewcommand{\thesection}{S\arabic{section}}

\title{Supplement to ``The least prime in an arithmetic progression: the exponent $278/125$ from a half-plane zero-free
region, a zero-density estimate and a formalised sieve''}
\author{Ritchie Comley}
\date{Draft, 28 October 2026}

\begin{document}
\maketitle

This supplement holds the proofs of four standard results that the paper uses, in the form in which the Lean library
proves them (Sections~\ref{app:prime-sum} to~\ref{app:rough}), and the details of the paper's account of the
formalisation (Section~\ref{sup:lean}) and of its computations around the exponent (Section~\ref{sup:computations}).
Numbers such as Lemma~\ref{M-lem:prime-sum} and Section~\ref{M-sec:certificate} refer to the paper, and those that
begin with S to this supplement. Its notation is the paper's.

\section{Proof of Lemma~\ref{M-lem:prime-sum}}\label{app:prime-sum}

\begin{proof}
Throughout, $Q \ge 4$, so that $\log Q > 1$, and $L_\chi$ stands for $L(\cdot, \chi)$, which is entire since $\chi$ is
not trivial. Put $\sigma_* = 11/12 + \delta$ and $\tau = (1 + H)^{-J}$.

\emph{The weight.} Let $f_Q(u) = f_0(u) \log Q/\log(uQ)$ where $f_0(u) \ne 0$, which needs $u \ge 1/2$ and so
$uQ \ge 2$, and $f_Q(u) = 0$ elsewhere; and let $w_y(u) = u^{-iy} f_Q(u)$. If
$f_0(n/Q) \ne 0$ then $n \ge Q/2 \ge 2$ and $f_Q(n/Q) = f_0(n/Q) \log Q/\log n$, so that
$\Lambda(n) f_Q(n/Q) = \log Q \cdot (\Lambda(n)/\log n) f_0(n/Q)$. Since $\Lambda(p)/\log p = 1$,
\begin{gather*}
  \sum_n \chi(n) \Lambda(n) w_y(n/Q) = \log Q \, \bigl(S_Q(\chi, y) + R\bigr), \\
  R = \sum_{\substack{p^k \le 2Q \\ k \ge 2}} \frac{\chi(p^k)}{k} \Bigl(\frac{p^k}{Q}\Bigr)^{-iy} f_0\Bigl(\frac{p^k}{Q}\Bigr).
\end{gather*}
Each term of $R$ is at most $C = \max |f_0|$ in absolute value, and a proper prime power $p^k \le 2Q$ has
$p \le \sqrt{2Q}$ and $k \le \log_2(2Q)$, so $|R| \le C(\sqrt{2Q} + 1)(\log_2(2Q) + 1)$. It remains to prove
\begin{equation}\label{eq:Sw-bound}
  \Bigl|\sum_n \chi(n) \Lambda(n) w_y(n/Q)\Bigr| \le M_\delta(\theta) \, \lambda \, B \bigl(Q^{\theta + \delta} + Q^{\sigma_*} \tau\bigr).
\end{equation}

\emph{The Mellin transform.} Write $\mathcal M f(s) = \int_0^\infty u^{s - 1} f(u) \, du$. The substitution
$u = e^{-v}$ gives, with $c_Q = 1/\log Q \in (0, 1/\log 4]$,
\[
  \mathcal M f_Q(\sigma + it) = \int_{\mathbb R} e^{-itv} V(v) \, dv, \qquad
  V(v) = e^{-\sigma v} f_0(e^{-v}) (1 - c_Q v)^{-1}.
\]
The function $V$ vanishes for $|v| > \log 2$, and $c_Q v \le 1/2$ for $v \le \log 2$; the pole $v = \log Q \ge \log 4$ of
$(1 - c_Q v)^{-1}$ lies where $V$ vanishes, so $V$ is smooth. By Leibniz's
rule, with $|(d/dv)^i (1 - c_Q v)^{-1}| = i! \, c_Q^i |1 - c_Q v|^{-1 - i} \le i! \, 2^{i + 1}$ on $|v| \le \log 2$,
the $j$-th derivative of $V$ is bounded by a constant $V_j$ that depends only on $f_0$ and $j$, uniformly in
$0 \le \sigma \le 2$ and $Q \ge 4$. Integration by parts $j$ times gives
$|t|^j |\mathcal M f_Q(\sigma + it)| \le 2 \log 2 \cdot V_j$, and with $j = 0$ and $j = J + 3$,
\[
  (1 + |t|)^{J + 3} |\mathcal M f_Q(\sigma + it)| \le B \qquad (0 \le \sigma \le 2,\ t \in \mathbb R,\ Q \ge 4),
\]
where $B = 2^{J + 4} \log 2 \, (V_0 + V_{J + 3})$ depends only on $f_0$ and $J$. Since $f_Q$ is continuous on
$(0, \infty)$ and vanishes near $0$ and near $\infty$, $\mathcal M f_Q$ is entire. Finally
$\mathcal M w_y(s) = \mathcal M f_Q(s - iy)$, so that
\begin{equation}\label{eq:mellin-decay}
  (1 + |t - y|)^{J + 3} |\mathcal M w_y(\sigma + it)| \le B \qquad (0 \le \sigma \le 2,\ t \in \mathbb R).
\end{equation}

\emph{Perron's formula.} For $\operatorname{Re} s > 1$, $-L_\chi'/L_\chi(s) = \sum_n \chi(n) \Lambda(n) n^{-s}$.
Mellin inversion on the line $\operatorname{Re} s = 2$ at each point $n/Q$, with the exchange of sum and integral
justified by \eqref{eq:mellin-decay}, gives
\begin{gather*}
  \sum_n \chi(n) \Lambda(n) w_y(n/Q) = \frac1{2\pi} \int_{\mathbb R} F(2 + it) \, dt, \\
  F(s) = -\frac{L_\chi'}{L_\chi}(s) \, \mathcal M w_y(s) \, Q^s.
\end{gather*}

\emph{The logarithmic derivative.} For $1/2 \le \theta_0 < \theta' < 3/2$ there is $\kappa(\theta_0, \theta')$ with the
following property, for every $q$, every non-trivial $\chi$ mod $q$ and every $t \in \mathbb R$: if $L_\chi$ has no
zero in the open disc $D_t$ of centre $3/2 + it$ and radius $\rho = 3/2 - \theta_0$, then
\begin{equation}\label{eq:logderiv}
  \Bigl|\frac{L_\chi'}{L_\chi}(\sigma + it)\Bigr| \le \kappa(\theta_0, \theta') \bigl(\log(q + 2) + \log(|t| + 4)\bigr)
  \qquad (\sigma \ge \theta').
\end{equation}
For $\sigma \ge 3/2$ this holds as soon as $\kappa \ge -\zeta'/\zeta(3/2)$, by the Dirichlet series of
$L_\chi'/L_\chi$. For $\theta' \le \sigma < 3/2$ it follows from a lemma of Landau's type: if $f$ is analytic near the
closed unit disc, $f(0) = 1$, $|f| \le B_0$ on $|z| \le R$ with $B_0 > 1$, $f$ has no zero in $|z| < 1$ and
$0 < r' < r < R' < R < 1$, then
\[
  |f'/f(z)| \le \Bigl(\frac{16 r^2}{(r - r')^3} + \frac1{(R^2/R' - R') \log(R/R')}\Bigr) \log B_0 \qquad (|z| \le r').
\]
(The library uses the form of this lemma in PrimeNumberTheoremAnd \cite{pnt-plus}, \lean{FinalBound}, which combines a
Blaschke product, the Borel--Carath\'eodory inequality and Jensen's formula, and allows zeros.) Apply it to
$f(z) = L_\chi(3/2 + it + \rho z)/L_\chi(3/2 + it)$ at $z = (\sigma - 3/2)/\rho$, with $r' = (3/2 - \theta')/\rho$ and
with $r$, $R'$ and $R$ a quarter, a half and three quarters of the way from $r'$ to $1$; then
$f'/f(z) = \rho \, L_\chi'/L_\chi(\sigma + it)$. The bound $B_0$ comes from two estimates. First,
$|L_\chi(s)| \le 3(q + 1)(|\operatorname{Im} s| + 2)$ for $1/2 \le \operatorname{Re} s \le 3$, which follows from the
partial-summation bound $|L_\chi(s)| \le (q + 1)|s|/\operatorname{Re} s$ for $\operatorname{Re} s > 0$. Second,
$|L_\chi(3/2 + it)|^{-1} \le \zeta(3/2)$, from $1/L_\chi(s) = \sum \mu(n) \chi(n) n^{-s}$. Together they give
$|f| \le B_0 = 2 + 6(q + 1)\zeta(3/2)(|t| + 4)$ on $|z| \le R$, and $\log B_0 \ll \log(q + 2) + \log(|t| + 4)$. The
constant $\kappa(\theta_0, \theta')$ is explicit and of order $(\theta' - \theta_0)^{-3}$.

The disc $D_t$ lies in $\operatorname{Re} s > \theta_0$ and, since $\rho \le 1$, in $|\operatorname{Im} s - t| < 1$. So
\eqref{eq:logderiv} applies in two cases:
\begin{enumerate}
\item[(i)] with $\theta_0 = 11/12$ and every $t$, by Input~\ref{M-inp:half-plane};
\item[(ii)] with $\theta_0 = \theta$ and every $t$ with $|t - y| \le H$, by the hypothesis of the lemma.
\end{enumerate}
Since $|t| + 4 \le (|y| + 4)(1 + |t - y|)$ and $\log(1 + u) \le u$, we have
$\log(q + 2) + \log(|t| + 4) \le \lambda(1 + |t - y|)$. With \eqref{eq:mellin-decay} this gives, in either case,
\begin{equation}\label{eq:F-bound}
  |F(\sigma + it)| \le \kappa(\theta_0, \theta') \, \lambda \, B \, Q^\sigma (1 + |t - y|)^{-J - 2} \qquad (\theta' \le \sigma \le 2).
\end{equation}

\emph{The contour.} By Input~\ref{M-inp:half-plane}, $F$ is holomorphic on $\sigma_* \le \operatorname{Re} s \le 2$, and by
\eqref{eq:F-bound} in case (i), with $\theta' = \sigma_*$, it is $O((1 + |t - y|)^{-2})$ there uniformly. So Cauchy's
theorem on the rectangles $[\sigma_*, 2] \times [-Y, Y]$, with $Y \to \infty$, moves the line of integration from $2$ to
$\sigma_*$. Split the line $\operatorname{Re} s = \sigma_*$ at the heights $y \pm H$. By the hypothesis of the lemma, $F$
is holomorphic on the rectangle $\theta + \delta \le \operatorname{Re} s \le \sigma_*$, $|\operatorname{Im} s - y| \le H$,
so by Cauchy's theorem the integral over $|t - y| \le H$ on the line $\sigma_*$ equals the integral over the same
heights on the line $\theta + \delta$ plus the integrals over the two horizontal sides at the heights $y \pm H$.

\emph{The pieces.} Put $\kappa_1 = \kappa(11/12, \sigma_*)$ and $\kappa_2 = \kappa(\theta, \theta + \delta)$.
\begin{itemize}
\item On the line $\sigma_*$ with $|t - y| \ge H$: by \eqref{eq:F-bound} in case (i) and
  $(1 + u)^{-J - 2} \le (1 + H)^{-J} (1 + u^2)^{-1}$ for $u \ge H$,
  $|F| \le \kappa_1 \lambda B Q^{\sigma_*} \tau/(1 + (t - y)^2)$. Each of the two tails contributes at most
  $\pi \kappa_1 \lambda B Q^{\sigma_*} \tau$.
\item On the line $\theta + \delta$ with $|t - y| \le H$: by \eqref{eq:F-bound} in case (ii),
  $|F| \le \kappa_2 \lambda B Q^{\theta + \delta}/(1 + (t - y)^2)$, so this part contributes at most
  $\pi \kappa_2 \lambda B Q^{\theta + \delta}$.
\item On each horizontal side, of length $11/12 - \theta \le 1$: by \eqref{eq:F-bound} in case (ii), with
  $Q^\sigma \le Q^{\sigma_*}$, $|F| \le \kappa_2 \lambda B Q^{\sigma_*} (1 + H)^{-J - 2} \le \kappa_2 \lambda B Q^{\sigma_*} \tau$.
\end{itemize}
Adding the pieces and dividing by $2\pi$, the sum in \eqref{eq:Sw-bound} is at most
\[
  \kappa_1 \lambda B Q^{\sigma_*} \tau + \tfrac12 \kappa_2 \lambda B Q^{\theta + \delta} + \tfrac1\pi \kappa_2 \lambda B Q^{\sigma_*} \tau
  \le (\kappa_1 + \kappa_2) \, \lambda \, B \bigl(Q^{\theta + \delta} + Q^{\sigma_*} \tau\bigr),
\]
which is \eqref{eq:Sw-bound} with $M_\delta(\theta) = \kappa(\theta, \theta + \delta) + \kappa(11/12, 11/12 + \delta)$.
\end{proof}

\section{Proof of Lemma~\ref{M-lem:family}}\label{app:family}

\begin{proof}
Let the family have the scale $P$, the set $S$, the coefficients $a_m$, the primes $p_1, \dots, p_K$ at the scales
$Q_i$, and the comparisons $C_1, \dots, C_J$ with forms $\ell_j$ and pivots $\pi(j)$. Let $\mathcal T$ be the set of the
$t = (m, p_1, \dots, p_K)$ with $m \in S$ and each $p_i$ a prime at most $2Q_i$; for $t \in \mathcal T$ put
$w(t) = a_m \prod_i f_0(p_i/Q_i)$ and $n(t) = m p_1 \cdots p_K$, and for a function $v$ on $\mathcal T$ put
\[
  \mathcal D(v) = \sum_{t \in \mathcal T} v(t) \, \Delta(n(t)), \qquad
  \Delta(n) = \mathbf 1[n \equiv a] - \frac{\mathbf 1[(n, q) = 1]}{\varphi(q)}.
\]
The sum of the lemma is $\mathcal D(w \prod_j I_j)$, with $I_j(t) = \mathbf 1[C_j \text{ holds at } t]$. Let $F_0 = \max f_0$ and
$\eta_0 = 10^{-5}$.

\emph{Orthogonality.} For $a$ prime to $q$,
$\Delta(n) = \varphi(q)^{-1} \sum_{\chi \in \mathcal X^*} \overline\chi(a) \chi(n)$, so for every $v$
\begin{equation}\label{eq:orth}
  |\mathcal D(v)| \le \frac1{\varphi(q)} \sum_{\chi \in \mathcal X^*} \Bigl|\sum_{t \in \mathcal T} v(t) \chi(n(t))\Bigr|.
\end{equation}

\emph{A smooth family.} Let $g_1, \dots, g_J$ be continuous and integrable, with integrable Fourier transforms
$\hat g_j(\xi) = \int g_j(u) e^{-2\pi i u \xi} \, du$, and let $T_0 \ge 1$ with
$2\pi \sum_j |e_{ji}| \, T_0 \le q^{\eta_0}$ for every $i$. We show that for $q$ larger than the $q_0$ of
Theorem~\ref{M-thm:type2} (with $f_0$, $K_0$, $e_1$ and $\eta_0$),
\begin{multline}\label{eq:smooth-family}
  \sum_{\chi \in \mathcal X^*} \Bigl|\sum_{t \in \mathcal T} w(t) \prod_j g_j(\ell_j(t)) \, \chi(n(t))\Bigr|
  \le q^{L - 1/50000} \mathcal N \prod_j \|\hat g_j\|_1 \\
  + \varphi(q) \, \mathfrak B \sum_j \int_{|\xi| > T_0} |\hat g_j| \prod_{l \ne j} \|\hat g_l\|_1,
\end{multline}
where
\[
  \mathcal N = \Bigl(\frac1P \sum_{m \in S} |a_m|^2\Bigr)^{1/2} \le A_0, \qquad
  \mathfrak B = \sum_m |a_m| \prod_i \sum_{p \le 2Q_i} f_0(p/Q_i).
\]
By Fourier inversion, and since the phase factorises over the variables,
\begin{gather*}
  \prod_j g_j(\ell_j) = \int_{\mathbb R^J} \prod_j \hat g_j(\xi_j) \, e^{2\pi i \sum_j \xi_j \ell_j} \, d\xi, \\
  e^{2\pi i \sum_j \xi_j \ell_j(t)} = e^{2\pi i \sum_j \xi_j \lambda_j(m)} \prod_i p_i^{2\pi i H_i}, \qquad
  H_i = \sum_j \xi_j e_{ji}.
\end{gather*}
So the inner sum in \eqref{eq:smooth-family} is
$\int \prod_j \hat g_j(\xi_j) \, W_\xi(\chi) \prod_i Q_i^{2\pi i H_i} S_{Q_i}(\chi, h_i) \, d\xi$, where
$W_\xi(\chi) = \sum_m a_m e^{2\pi i \sum_j \xi_j \lambda_j(m)} \chi(m)$ and $h_i = -2\pi H_i$. Take absolute values,
sum over $\chi$, and bound the integrand for each $\xi$. If $|\xi_j| \le T_0$ for every $j$, then $|h_i| \le q^{\eta_0}$,
and Theorem~\ref{M-thm:type2}, with the coefficients $a_m e^{2\pi i \sum_j \xi_j \lambda_j(m)}$, which have the moduli
$|a_m|$, and the heights $h_i$, gives
$\sum_\chi |W_\xi(\chi)| \prod_i |S_{Q_i}(\chi, h_i)| \le q^{L - 1/50000} \mathcal N$. Otherwise some $|\xi_j| > T_0$, and the
trivial bound $\varphi(q) \mathfrak B$ holds. Integrating gives \eqref{eq:smooth-family}.

\emph{The cut-off.} Let $g$ be smooth, with $g = 0$ on $(-\infty, 0]$, $g = 1$ on $[1, \infty)$ and $0 \le g \le 1$, and
let $c_k = \int |g^{(k)}|$. For $0 < \nu \le 2\Lambda$ put $G_{\nu, \Lambda}(u) = g(u/\nu) - g((u - 2\Lambda)/\nu)$.
It is smooth, $0 \le G_{\nu, \Lambda} \le 1$, it vanishes outside $[0, 2\Lambda + \nu]$, and
$G_{\nu, \Lambda}(u) = g(u/\nu)$ for $u \le 2\Lambda$. So if $u \le \Lambda$ and a condition holds when $u > 0$ and
fails when $u < 0$,
\begin{equation}\label{eq:cutoff}
  \bigl|\mathbf 1[\text{condition}] - G_{\nu, \Lambda}(u)\bigr| \le G_{\nu, \nu}(u + \nu),
\end{equation}
since the left side vanishes unless $0 \le u < \nu$, where $G_{\nu, \nu}(u + \nu) = 1$. Integration by parts
gives $(2\pi |\xi|)^k |\hat G_{\nu, \Lambda}(\xi)| \le \|G_{\nu, \Lambda}^{(k)}\|_1 \le 2\nu^{1 - k} c_k$, and
$|\hat G_{\nu, \Lambda}| \le 2\Lambda + \nu$. Hence
\begin{gather*}
  \|\hat G_{\nu, \Lambda}\|_1 \le \pi(2\Lambda + \nu) + \frac{c_2}{2\pi\nu}, \\
  \int_{|\xi| > T_0} |\hat G_{\nu, \Lambda}| \le \frac{4\pi c_k}{(2\pi)^k} \, \nu^{1 - k} T_0^{2 - k}
  \qquad (k \ge 2,\ T_0 \ge 1).
\end{gather*}

\emph{Sharp against smooth.} Let $\Lambda = \Lambda_0 \log q + 1$, so that $\ell_j \le \Lambda$ on $\mathcal T$, and put
$G_j = G_{\nu, \Lambda}(\ell_j)$ and $M_j = G_{\nu, \nu}(\ell_j + \nu)$. All these values lie in $[0, 1]$, so
by \eqref{eq:cutoff} $|\prod_j I_j - \prod_j G_j| \le \sum_j |I_j - G_j| \le \sum_j M_j$. Since
$|\Delta(n)| \le \Delta(n) + 2 \cdot \mathbf 1[(n, q) = 1]/\varphi(q)$ and $f_0 \ge 0$,
\begin{equation}\label{eq:sharp-smooth}
  \Bigl|\mathcal D\Bigl(w \prod_j I_j\Bigr)\Bigr| \le \Bigl|\mathcal D\Bigl(w \prod_j G_j\Bigr)\Bigr|
  + \sum_j \Bigl(|\mathcal D(|w| M_j)| + \frac2{\varphi(q)} \sum_{t \in \mathcal T} |w(t)| M_j(t)\Bigr).
\end{equation}
The first term and each $\mathcal D(|w| M_j)$ are smooth families: the first with $g_j = G_{\nu, \Lambda}$, the second with
the coefficients $|a_m|$ and the single function $G_{\nu, \nu}$ at the form $\ell_j + \nu$. Both are bounded by
\eqref{eq:orth} and \eqref{eq:smooth-family}. A trivial count would not do for them, since it would lose the factor $1/\varphi(q)$ that the condition
$n \equiv a$ gives. The last sum, which carries that factor already, is counted trivially, by the pivot. Fix $m$ and the primes
other than $p_{\pi(j)}$; then $\ell_j = c + s \log p_{\pi(j)}$ with $s = \pm 1$, and $M_j \ne 0$ forces
$-\nu \le \ell_j \le 2\nu$, which confines $p_{\pi(j)}$ to an interval $[Y, e^{3\nu} Y]$. At most
$(e^{3\nu} - 1) 2Q_{\pi(j)} + 1$ integers up to $2Q_{\pi(j)}$ lie in it, so
\[
  \sum_{t \in \mathcal T} |w(t)| M_j(t)
  \le F_0 \bigl((e^{3\nu} - 1) 2Q_{\pi(j)} + 1\bigr) \sum_m |a_m| \prod_{i \ne \pi(j)} \sum_{p \le 2Q_i} f_0(p/Q_i).
\]

\emph{The parameters.} Put $\varepsilon_2 = 1/(400000(J_0 + 1))$, $\nu = q^{-\varepsilon_2}$, $T_0 = q^{1/200000}$,
$k = 1200003$ and $\delta = \min(\varepsilon_2, e_1)/2$. For large $q$:
\begin{itemize}
\item $2\pi J_0 E_0 T_0 \le q^{\eta_0}$, so \eqref{eq:smooth-family} applies, as $\sum_j |e_{ji}| \le J_0 E_0$;
\item $\pi(2\Lambda + 1) \le q^{\varepsilon_2}$, so $\|\hat G_{\nu, \Lambda}\|_1$ and $\|\hat G_{\nu, \nu}\|_1$
  are at most $C q^{\varepsilon_2}$ with $C = 1 + c_2/(2\pi)$, and a product of at most $J_0$ of them is at most
  $C^{J_0 + 1} q^{1/400000}$;
\item $\nu^{1 - k} T_0^{2 - k} = q^{\varepsilon_2(k - 1) - (k - 2)/200000} \le q^{-3}$, which bounds the tails.
\end{itemize}
The constants come from $\eta_0 = 10^{-5}$: the norms lose at most $q^{\eta_0/4}$, $T_0 = q^{\eta_0/2}$ leaves room
for the factor $2\pi J_0 E_0$ below $q^{\eta_0}$, and $k = 1200003$ is the least $k$ with
$(k - 1)\eta_0/4 - (k - 2)\eta_0/2 \le -3$, so that the tails are at most $q^{-3}$ whatever $J_0$.
Moreover $\mathfrak B \le P A_0 (3F_0)^K \prod_i Q_i \le A_0 (3F_0 + 1)^{K_0} 2^{K_0} q^L$, since at most $3Q$ integers
lie in $[1, 2Q]$ for $Q \ge 1$ and $P \prod_i Q_i \le 2^{K_0} x$. Finally $e^{3\nu} - 1 \le 6\nu$ for
$\nu \le 1/3$ and $Q_{\pi(j)} \ge q^{e_1}$ give
$(12\nu Q_{\pi(j)} + 1) \prod_{i \ne \pi(j)} Q_i \le (12\nu + q^{-e_1}) \prod_i Q_i$. Putting these into
\eqref{eq:sharp-smooth}, with \eqref{eq:orth} and \eqref{eq:smooth-family}, and using $\varphi(q) \le q$, the sum of the
lemma is at most
\[
  \frac{A_0 q^L}{\varphi(q)} \bigl(a_1 q^{-7/400000} + a_2 q^{-1} + a_3(12 q^{-\varepsilon_2} + q^{-e_1})\bigr),
\]
where $a_1$, $a_2$ and $a_3$ depend only on $K_0$, $J_0$ and $F_0$. The three terms come from the smooth families, from
the tails and from the pivot counts. Each power of $q$ in the bracket is at most $q^{-2\delta}$, so for large $q$ the sum
is at most $A_0 q^{L - \delta}/\varphi(q)$. Here $\delta$ depends only on $J_0$ and $e_1$, and Theorem~\ref{M-thm:type2}
was used once for the smoothed family and once for each near-boundary family.
\end{proof}

\section{Proof of Proposition~\ref{M-prop:harman}}\label{app:harman}

\begin{proof}
Write $\mathcal P[y, z)$ for the set of primes in $[y, z)$, and say that $n$ is \emph{$z$-rough} if it has no prime
factor below $z$. For $\mathcal C = \mathcal A$ or $\mathcal B$ and $d$ whose prime factors are all at least $z$,
$S(\mathcal C_d, z)$ is the number of $z$-rough multiples of $d$ in $\mathcal C$. Throughout, $x = q^L$, and the
integers of $(x, 2x]$ are those of $(N_1, N_2]$ with $N_1 = \lfloor x \rfloor$ and $N_2 = \lfloor 2x \rfloor$. The tuples $(r_1, \dots, r_j)$ are ordered, and $r = r_1 \cdots r_j$.

\emph{The constants.} They are chosen before $q$ and depend only on $L$, $u_1$, $u_2$, $j$, $J_0$, $\Lambda_1$, $a_0$,
$b_0$ and $\varepsilon_H$. Let $c = (L - 1 - b_0)/4 > 0$. Let $C_0 \ge 1$ be an absolute constant with
$\sum_{p \in \mathcal P[y, y^2)} 1/p \le C_0$ for every $y \ge 1$; it exists by Mertens' theorem (the library
takes it from the prime number theorem). Put
$K_1 = (\log((L + 1)/a_0) + 1)^j$ and
\[
  \zeta_k = e^{2C_0} \cdot \frac{2 C_0^k}{k!} \cdot \exp\Bigl(e^{2C_0} \frac{2 C_0^{k + 1}}{(k + 1)!}\Bigr).
\]
Since $(2k + 4)^2 \zeta_k \to 0$, there is an even $k \ge 2$ with $2C_0^2 \le (k + 1)(k + 2)$, $c < (2k + 4) a_0$
and $(2k + 4)^2 \zeta_k \le \varepsilon_H c^2/(4e(K_1 + 1) a_0)$. Put $\varepsilon_w = c/(2k + 4)$, so
that $0 < \varepsilon_w < a_0$ and $(2k + 7) \varepsilon_w \le 2c$, and $S_0 = \lceil L/\varepsilon_w \rceil$. At $q$,
put $w = \lceil q^{\varepsilon_w} \rceil$, $z_0 = \lceil q^{a_0} \rceil$ and $X' = \lfloor q^{b_0} \rfloor$. Then
$w \le z_0$, and by hypothesis every tuple in the sum has $r_i \ge z_0$ and $r \le X'$.

\emph{The tree identity.} Let $r$ have all its prime factors at least $z_0$, and let $r \le X'$. For a set $P$ of primes
write $d_P = r \prod_{p \in P} p$, and let $\rho_z(n) = 1$ if $n$ is $z$-rough and $\rho_z(n) = 0$ otherwise. Then for
every $n$,
\begin{multline}\label{eq:tree}
  \mathbf 1[r \mid n] \, \rho_{z_0}(n)
  = \sum_{P \subseteq \mathcal P[w, z_0)} (-1)^{|P|} \, \mathbf 1[d_P \le X'] \, \mathbf 1[d_P \mid n] \, \rho_w(n) \\
  + \sum_{p \in \mathcal P[w, z_0)} \ \sum_{P \subseteq \mathcal P(p, z_0)} (-1)^{|P| + 1} \,
    \mathbf 1[d_P \le X' < p \, d_P] \, \mathbf 1[p \, d_P \mid n] \, \rho_p(n),
\end{multline}
where $\mathcal P(p, z_0)$ is the set of primes in $(p, z_0)$. This is what iterating Buchstab's identity
$S(\mathcal C_m, z) = S(\mathcal C_m, w) - \sum_{w \le p < z} S(\mathcal C_{mp}, p)$ from the root $r$ at $z_0$ gives,
when a child $mp$ is stopped as soon as $mp > X'$: the first sum collects the \emph{leaves} $S(\mathcal C_d, w)$,
$d = d_P \le X'$, and the second the \emph{stops} $S(\mathcal C_W, p)$, $W = p \, d_P$. It is proved directly,
for each $n$. If $r \nmid n$ or $n$ is not $w$-rough, every term vanishes. Otherwise let $Q$ be the set of the primes of
$[w, z_0)$ that divide $n$; the left side is $\mathbf 1[Q = \emptyset]$. Since the primes of $r$ are at least $z_0$, a
leaf $P$ counts $n$ exactly when $P \subseteq Q$, and a stop $(p, P)$ exactly when $p = \min Q$ and
$P \subseteq Q \setminus \{p\}$. If $Q = \emptyset$, only the leaf $P = \emptyset$ counts $n$, with
$\mathbf 1[r \le X'] = 1$. If $Q \ne \emptyset$, let $p_0 = \min Q$; grouping $P \subseteq Q \setminus \{p_0\}$ with
$P \cup \{p_0\}$, and writing $R = d_P$, the right side is
\[
  \sum_{P \subseteq Q \setminus \{p_0\}} (-1)^{|P|}
  \bigl(\mathbf 1[R \le X'] - \mathbf 1[p_0 R \le X'] - \mathbf 1[R \le X' < p_0 R]\bigr) = 0.
\]
Multiplying \eqref{eq:tree} by $\mathbf 1[n \equiv a] - \mathbf 1[(n, q) = 1]/\varphi(q)$ and summing over $n$ and over
the tuples, the sum of the proposition is the sum over the leaves of
$(-1)^{|P|}(S(\mathcal A_d, w) - S(\mathcal B_d, w)/\varphi(q))$ plus the sum over the stops of
$(-1)^{|P| + 1}(S(\mathcal A_W, p) - S(\mathcal B_W, p)/\varphi(q))$.

\emph{The stops.} A stop has $X' < W \le p X' < z_0 X'$, and $W$ is $p$-rough, since the primes of $r$ are at least
$z_0 > p$ and those of $P$ exceed $p$. So the $p$-rough multiples of $W$ in $(N_1, N_2]$ are the $n = Wm$ whose cofactor
$m$ is $p$-rough, and every such $m$ lies in $(P_1, P_2]$ with $P_1 = \lfloor N_1/(z_0 X') \rfloor$ and
$P_2 = \lfloor N_2/(X' + 1) \rfloor$. For large $q$, $q^{u_1 L} \le P_1$, since $5q^{u_1 L + a_0 + b_0} \le q^L$ by
$a_0 + b_0 < (1 - u_1)L$, and $P_2 \le q^{u_2 L}$, since $2q^{L - b_0} \le q^{u_2 L}$ by $(1 - u_2)L < b_0$. So the
cofactor $m$ lies in the window's range, and $W$ in its reflection. A stop has $|P| \le S_0$, since
$\prod P \ge w^{|P|} \ge q^{\varepsilon_w |P|}$ and $\prod P \le X' \le q^L$.

For each $s \le S_0$, the stops with $|P| = s$ form one sum of the kind that Lemma~\ref{M-lem:family-sharp} bounds, with
$m \in (P_1, P_2]$ and $a_m = 1$, whose primes are $r_1, \dots, r_j$, the elements $y_1 > \dots > y_s$ of $P$ and
$p$, and whose comparisons are the given ones on the $r_i$ and
\begin{gather*}
  y_1 > \dots > y_s > p, \quad y_1 < z_0, \quad p \ge w, \quad r y_1 \cdots y_s \le X' < r y_1 \cdots y_s \, p, \\
  p \le \min\{\ell \text{ prime} : \ell \mid m\}, \qquad N_1 < m r y_1 \cdots y_s p \le N_2 .
\end{gather*}
For $s = 0$ the condition $y_1 < z_0$ reads $p < z_0$. Each condition is of the form
$\lambda(m) + \sum_i e_i \log x_i > 0$ or $\ge 0$ with a pivot: the new ones have coefficients in $\{-1, 0, 1\}$,
and the given ones coefficients at most $1$ in absolute value. There are at most $J_0 + s + 7$ of them, and
$\lambda(m) \le (\Lambda_1 + L + 1) \log q$ for $m \le P_2$. Wherever they
all hold, every prime is at least $w \ge 2q^{\varepsilon_w/2}$ and $m r y_1 \cdots y_s p \le 2x$. So
Lemma~\ref{M-lem:family-sharp}, with $K_0 = j + S_0 + 2$, $J_0 + S_0 + 7$ comparisons, $e_1 = \varepsilon_w/2$, $E_0 = 1$
and $\Lambda_1 + L + 1$ in place of $\Lambda_1$, bounds each of these $S_0 + 1$ sums by $q^{L - \delta_T}/\varphi(q)$,
with $\delta_T > 0$ depending only on the constants above. For large $q$ the stops therefore total at most
$q^{L - \delta_1}/\varphi(q)$, with $\delta_1 = \delta_T/2$.

\emph{A leaf.} Let $d$ have all its prime factors at least $w$. We show
\begin{equation}\label{eq:leaf}
  \Bigl|S(\mathcal A_d, w) - \frac{S(\mathcal B_d, w)}{\varphi(q)}\Bigr|
  \le \frac{2 \zeta_k (N_2 - N_1)}{d \varphi(q) \log w} + 24 \, w^{2k + 7}.
\end{equation}
If $(d, q) > 1$ both counts vanish, so let $(d, q) = 1$. Sieve by the primes below $w$, in the ranges
$T_i = \mathcal P[w^{2^{-i - 1}}, w^{2^{-i}})$ for $0 \le i < J = 2 \lfloor \log_2 w \rfloor + 2$, which cover them, at the
even levels $k_i = k + 2i$, which grow with $i$ so that the errors $C_0^{k_i}/k_i!$ below have a sum that does not
depend on $J$. Brun's sieve \cite{brun}, in the following form, is used, for disjoint sets $T_i$ of
primes and even levels $k_i \ge 2$. Let $\mathcal C$ be a finite set of integers,
$X \ge 0$, and $g$ a function on the primes of $\bigcup T_i$ with $0 \le g \le 1$. Suppose that for every choice of
$s_i \subseteq T_i$, with $D = \prod_i \prod_{p \in s_i} p$, the count $\#\{n \in \mathcal C : D \mid n\}$ is
$X \prod_i \prod_{p \in s_i} g(p)$ within $R$. Then $\#\{n \in \mathcal C : (n, \prod_i \prod_{T_i} p) = 1\}$ is
$X \prod_i \prod_{p \in T_i} (1 - g(p))$ within
\[
  X \sum_i e_{k_i}(T_i) \prod_{l \ne i} F_{k_l}(T_l)
  + (J + 1) \, R \prod_i \sum_{m \le k_i} \binom{|T_i|}{m},
\]
where $e_m(T) = \sum_{s \subseteq T,\ |s| = m} \prod_{p \in s} g(p)$ and $F_k(T) = \sum_{m \le k} (-1)^m e_m(T)$.
Bonferroni's inequalities say that $F_{k - 1}(T) \le \prod_{p \in T}(1 - g(p)) \le F_k(T)$ for every even $k$, and the
same for any weights in $[0, 1]$ in place of $g$. Take them for each $n \in \mathcal C$ with the weights
$\mathbf 1[p \mid n]$, and use the elementary inequality
$\prod_i \mathrm{hi}_i - \sum_i (\mathrm{hi}_i - \mathrm{lo}_i) \prod_{l \ne i} \mathrm{hi}_l \le \prod_i \rho_i$ for
$\rho_i \ge 0$ and $\mathrm{lo}_i \le \rho_i \le \mathrm{hi}_i$ (the $\mathrm{lo}_i$, odd truncations, may be negative).
Summed over $n$, and expanded by the hypothesis, they place the count, up to $(J + 1)$ times
$R \prod_i \sum_{m \le k_i} \binom{|T_i|}{m}$, in the interval from $X \prod_i F_{k_i} - X \sum_i e_{k_i} \prod_{l \ne i}
F_{k_l}$ to $X \prod_i F_{k_i}$. The model $X \prod_i V_i$, with $V_i = \prod_{p \in T_i}(1 - g(p))$, lies in the same
interval: Bonferroni's inequalities for $g$ give $F_{k_i} - e_{k_i} = F_{k_i - 1} \le V_i \le F_{k_i}$, hence
$0 \le \prod_i F_{k_i} - \prod_i V_i \le \sum_i e_{k_i} \prod_{l \ne i} F_{k_l}$.

Take $g(p) = 1/p$ for $p \nmid q$ and $g(p) = 0$ for $p \mid q$. For $\mathcal A_d$, take
$\mathcal C = \{n \in (N_1, N_2] : n \equiv a,\ d \mid n\}$, $X = (N_2 - N_1)/(dq)$ and $R = 1$: if no prime of $D$
divides $q$, Lemma~\ref{M-lem:type1} with $M = dD$ gives the count within $1$, and otherwise the count and the model both
vanish. For $\mathcal B_d$, take $\mathcal C = \{n \in (N_1, N_2] : (n, q) = 1,\ d \mid n\}$, $X\varphi(q)$ and
$R = \varphi(q)$, and divide by $\varphi(q)$. Both main terms are then $XV$ with $V = \prod_{p < w,\ p \nmid q}(1 - 1/p)$,
so that
\begin{gather*}
  \Bigl|S(\mathcal A_d, w) - \frac{S(\mathcal B_d, w)}{\varphi(q)}\Bigr| \le 2(XG + \mathrm{rem}), \\
  G = \sum_i e_{k_i}(T_i) \prod_{l \ne i} F_{k_l}(T_l), \qquad
  \mathrm{rem} = (J + 1) \prod_i \sum_{m \le k_i} \binom{|T_i|}{m}.
\end{gather*}
Each $T_i$ lies in some $\mathcal P[y, y^2)$, so $\sum_{p \in T_i} g(p) \le C_0$ and, since $1 - t \ge e^{-2t}$ for
$0 \le t \le 1/2$, $V_i \ge e^{-2C_0}$. With $e_m \le C_0^m/m!$,
$F_{k_l} \le V_l + e_{k_l + 1} \le V_l (1 + e^{2C_0} C_0^{k_l + 1}/(k_l + 1)!)$ and
$\prod_{l \ne i} V_l \le e^{2C_0} V$, this gives
$G \le V e^{2C_0} \sum_i \frac{C_0^{k_i}}{k_i!} \prod_l \bigl(1 + e^{2C_0} \frac{C_0^{k_l + 1}}{(k_l + 1)!}\bigr)$. Since
$2C_0^2 \le (k + 1)(k + 2)$, the terms of $\sum_i C_0^{k + 2i}/(k + 2i)!$ decrease at least geometrically with ratio
$1/2$, so the sum is at most $2C_0^k/k!$, and with $\prod (1 + t_l) \le \exp \sum t_l$ we get $G \le \zeta_k V$. So
the difference is at most $2\zeta_k XV + 2\,\mathrm{rem}$. Next, $V \le q/(\varphi(q) \log w)$. Indeed
$V \prod_{p < w,\ p \mid q}(1 - 1/p) = \prod_{p < w}(1 - 1/p)$, the second product on the left is at least
$\prod_{p \mid q}(1 - 1/p) = \varphi(q)/q$, and $\prod_{p < w}(1 - 1/p) \le 1/\log w$, since
$\prod_{p < w}(1 - 1/p)^{-1}$ is the sum of $1/m$ over the integers $m \ge 1$ whose prime factors are all below $w$,
which is at least $\sum_{m < w} 1/m \ge \log w$. So $XV \le (N_2 - N_1)/(d \varphi(q) \log w)$. Finally
$\sum_{m \le k_i} \binom{|T_i|}{m} \le 2 \max(1, |T_i|)^{k_i}$ with $|T_i| \le w^{2^{-i}}$,
$\sum_i 2^{-i}(k + 2i) \le 2k + 4$, $J + 1 \le 3w$ and $2^J \le 4w^2$ give $\mathrm{rem} \le 12 w^{2k + 7}$.

\emph{The leaves together.} In \eqref{eq:leaf}, $N_2 - N_1 \le x + 1 \le 2x$ and $\log w \ge \varepsilon_w \log q$.
With $1 \le X'/d$ for the remainder, the leaves total at most
\[
  \Bigl(\frac{4 \zeta_k x}{\varphi(q) \varepsilon_w \log q} + 24 \, w^{2k + 7} X'\Bigr)
  \sum_{(r_1, \dots, r_j)} \sum_{P \subseteq \mathcal P[w, z_0)} \frac1{r \prod P},
\]
and the double sum is $\sum 1/r \cdot \prod_{w \le p < z_0}(1 + 1/p)$. Each $r_i$ is a prime in $[q^{a_0}, 2q^L]$, so
for large $q$, by Mertens' theorem (in the library, the prime number theorem), $\sum 1/r \le K_1$ and $\prod_{w \le p < z_0}(1 + 1/p) \le e a_0/\varepsilon_w$.
The main part is then at most
\[
  \frac{4 e K_1 a_0 \zeta_k}{\varepsilon_w^2} \cdot \frac{x}{\varphi(q) \log q}
  = \frac{4 e K_1 a_0 (2k + 4)^2 \zeta_k}{c^2} \cdot \frac{x}{\varphi(q) \log q}
  \le \frac{\varepsilon_H \, x}{\varphi(q) \log q},
\]
by the choice of $k$. Since $w \le 2q^{\varepsilon_w}$ and $(2k + 7)\varepsilon_w \le 2c$, the remainder part is at most
$24 \cdot 2^{2k + 7} K_1 e a_0 \varepsilon_w^{-1} q^{b_0 + 2c} \le q^{b_0 + 3c} = q^{L - 1 - c} \le q^{L - c}/\varphi(q)$
for large $q$, using $\varphi(q) \le q$.

\emph{Conclusion.} The sum of the proposition is at most
$\varepsilon_H x/(\varphi(q) \log q) + (q^{L - c} + q^{L - \delta_1})/\varphi(q)$, which for large $q$ is at most
$\varepsilon_H x/(\varphi(q) \log q) + q^{L - \delta}/\varphi(q)$ with $\delta = \min(c, \delta_1)/2$.
\end{proof}

\section{Proof of Lemma~\ref{M-lem:rough}}\label{app:rough}

\begin{proof}
Write $c_* = 0.565073$ and $\bar\omega = 0.5671433$, $u = \log 2y/\log z$, and say that an integer is $z$-rough if it has no
prime factor below $z$. Four facts are used.

\emph{(a) Prime sums.} By the prime number theorem with an error $o(t/\log t)$ (in the library, the bound
$\psi(t) - t = O(t^{11/12 + \varepsilon})$, which it derives from the case of $\zeta$ of Input~\ref{M-inp:half-plane};
the classical theorem would do), for every $\varepsilon_3 > 0$,
$|\theta(t) - t| \le \varepsilon_3 t/\log t$ for large $t$. With partial summation this gives, for large $a \le b$,
\[
  \sum_{a < p \le b} \frac1{p \log p} \le \frac1{\log a} - \frac1{\log b} + \frac{4\varepsilon_3}{\log^2 a},
\]
and, if $2 \log b \le \Lambda$,
\[
  \sum_{a < p \le b} \frac1{p(\Lambda - \log p)}
  \le \frac{1 + \varepsilon_3/\log a}{\Lambda} \log \frac{\log b \, (\Lambda - \log a)}{\log a \, (\Lambda - \log b)}
  + \frac{6\varepsilon_3}{\Lambda \log^2 a};
\]
and, for large $y$, there are at most $(1 + \varepsilon_3) y/\log 2y$ primes in $(y, 2y]$.

\emph{(b) Buchstab's inequality.} For $y \ge 2$,
\[
  \Phi(y, z) \le \#\{p : y < p \le 2y\} + \sum_{z \le p \le \sqrt{2y}} \Phi(y/p, p),
\]
since a $z$-rough integer $m \in (y, 2y]$ that is not prime has its least prime factor $p$ in $[z, \sqrt{2y}]$, and
$m/p$ is a $p$-rough integer in $(y/p, 2y/p]$; the map $m \mapsto (p, m/p)$ is injective. Also $\Phi(y, z) = 0$ if
$y \ge 1$ and $2y < z$, since every integer of $(y, 2y]$ is then at least $2$ and below $z$.

\emph{(c) Buchstab's function on $[2, 3]$.} For $u > 1$, $1 + \log(u - 1) \le \bar\omega u$: the inequality $\log t \le t - 1$ at
$t = \bar\omega(u - 1)$ gives $1 + \log(u - 1) \le \bar\omega(u - 1) - \log \bar\omega \le \bar\omega u$, since $-\log \bar\omega \le \bar\omega$, that is $\bar\omega e^{\bar\omega} \ge 1$, which
holds because $e^{\bar\omega} \ge \sum_{m \le 12} \bar\omega^m/m! > 1/\bar\omega$. So $\omega^* \le \bar\omega$ on $[2, 3]$, and since $c_* < \bar\omega$,
$\omega^* \le \bar\omega$ on $[2, \infty)$; and $\omega^* \le c_*$ on $[3, \infty)$, since $\omega^*(3) = (1 + \log 2)/3 < c_*$.
On $[2, 3]$, $\omega^*(u) = (1 + \log(u - 1))/u$ increases up to $2.76$ and decreases from $2.77$ on. For
$2 \le u \le u' \le 2.76$, $\omega^*(u) \le \omega^*(u')$ is $(u' - u)(1 + \log(u - 1)) \le u \log((u' - 1)/(u - 1))$:
the right side is at least $u(u' - u)/(u' - 1)$, by $\log t \ge 1 - 1/t$, and the left at most $(u' - u)u/1.76$, by
$\log t \le t - 1$ at $t = (u - 1)/1.76$ and $\log 1.76 \le 1/1.76$. For $2.77 \le u \le u'$ the reverse inequality
follows in the same way from $\log((u' - 1)/(u - 1)) \le (u' - u)/(u - 1)$ and
$\log(u - 1) \ge \log 1.77 \ge 1/1.77 \ge 1/(u - 1)$. The two constants come from the exponential's series. (This
proof does not use the monotonicity; the bound $\Omega$ of Section~\ref{M-sec:certificate} does.)

\emph{(d) A delay inequality.} Let $M(u) = \log(u - 1)$ for $2 \le u \le 3$ and
$M(u) = \log 2 + \bar\omega \min(u - 3, 1) + c_* \max(u - 4, 0)$ for $u > 3$, which is $\int_1^{u - 1} \omega^*(t)\, dt$ with
$\omega^*$ replaced by $\bar\omega$ on $[2, 3]$. Then $1 + M(u) \le u \omega^*(u)$ for $u \ge 2$. For $u \le 3$ it is an
equality. For $u > 3$, $\omega^*(u) = c_*$, and $1 + M(u) - c_* u$ is affine on $[3, 4]$, with slope $\bar\omega - c_* > 0$, and
constant on $[4, \infty)$; so it suffices that it is at most $0$ at $u = 4$, that is $1 + \log 2 + \bar\omega \le 4c_*$, which
holds since $\log 2 < 0.6931471808$ and $2.2602904808 < 2.260292 = 4c_*$.

We prove by induction on the integer $k \ge 0$ the statement $\mathrm C(k)$: for every $\varepsilon > 0$ there is
$z_0$ such that $\Phi(y, z) \le (\omega^*(u) + \varepsilon) y/\log z$ for all $z \ge z_0$ and all $y \ge 1$ with
$u \le k$. $\mathrm C(0)$ holds trivially, since $u > 0$, and the lemma is $\mathrm C(\lceil U \rceil)$. Assume
$\mathrm C(k)$, let $0 < \varepsilon \le 1$ (which suffices, since the bound for one $\varepsilon$ implies it for every
larger one) and $\varepsilon_3 = \varepsilon/100$, let $z$ be large in terms of
$\varepsilon$ and of the $z_0$ of $\mathrm C(k)$ for $\varepsilon/4$, and let $y \ge 1$ with
$\Lambda = \log 2y = u \log z \le (k + 1) \log z$. If $2y < z$, then $\Phi(y, z) = 0$. Otherwise $u \ge 1$, $y \ge z/2$
is large, and (a) and (b) give
\[
  \Phi(y, z) \le \Bigl(1 + \frac\varepsilon4\Bigr) \frac y\Lambda + \sum_{z \le p \le \sqrt{2y}} \Phi(y/p, p).
\]
If $u < 2$, then $\sqrt{2y} < z$, the sum is empty, and
$\Phi(y, z) \le (1 + \varepsilon/4) y/(u \log z) \le (\omega^*(u) + \varepsilon) y/\log z$. Let $u \ge 2$. A prime
$p \in [z, \sqrt{2y}]$ has $y/p \ge 1$ and
$\log(2y/p) = \Lambda - \log p \le k \log p$, since $\Lambda \le (k + 1) \log p$. So $\mathrm C(k)$ gives
$\Phi(y/p, p) \le (\omega^*(u_p) + \varepsilon/4) \, y/(p \log p)$ with $u_p = (\Lambda - \log p)/\log p$. The term
$p = z$, if $z$ is prime, contributes at most $2y/(z \log z)$, since $\omega^* \le 1$. By (a), the parts
$\varepsilon/4$ of the other terms contribute at most $(\varepsilon/4)(1/\log z + 4\varepsilon_3/\log^2 z) \, y$. For
the parts $\omega^*(u_p)$: if $p \le (2y)^{1/4}$ then $u_p \ge 3$ and $\omega^*(u_p) \le c_*$; if
$(2y)^{1/4} < p \le (2y)^{1/3}$ then $2 \le u_p < 3$ and $\omega^*(u_p) \le \bar\omega$; and if $(2y)^{1/3} < p \le (2y)^{1/2}$ then
$1 \le u_p < 2$ and $\omega^*(u_p)/(p \log p) = 1/(p(\Lambda - \log p))$, by the definition of $\omega^*$ on
$[1, 2)$. By (a), applied on $(z, \sqrt{2y}]$ if
$u \le 3$, on $(z, (2y)^{1/3}]$ and $((2y)^{1/3}, \sqrt{2y}]$ if $3 < u \le 4$, and also with the first range split at
$(2y)^{1/4}$ if $u > 4$,
\[
  \sum_{z < p \le \sqrt{2y}} \frac{\omega^*(u_p)}{p \log p} \le \frac{M(u) + 12\varepsilon_3}\Lambda
  + \frac{8\varepsilon_3}{\log z}.
\]
Collecting the terms, with $\Lambda = u \log z$,
\[
  \Phi(y, z) \le \frac y{\log z} \Bigl(\frac{1 + M(u)}u + \frac{\varepsilon/4 + 12\varepsilon_3}u + \frac2z
  + 8\varepsilon_3 + \frac\varepsilon4 + \frac{\varepsilon \varepsilon_3}{\log z}\Bigr).
\]
By (d) the first term is at most $\omega^*(u)$, and since $u \ge 1$ and $z$ is large the others add up to at most
$0.96\varepsilon$. So $\Phi(y, z) \le (\omega^*(u) + \varepsilon) y/\log z$, which proves $\mathrm C(k + 1)$.
\end{proof}

% The entries are the paper's; main.tex records how they were checked.
\section{The formalisation in detail}\label{sup:lean}

This section completes Section~\ref{M-sec:lean} of the paper.

\subsection*{Chen, Gupta and Li's paper}

For Chen, Gupta and Li's paper \cite{cgl} the library formalises its Proposition~3.1, in the form for one point per
character of Section~\ref{M-sec:inputs} of the paper (\lean{LinnikCGLLV.CGLProp31}, proved as
\lean{LinnikCGLHB101.prop31}), with the results it rests on (their Propositions~4.6, 6.1, 8.1 and~11.1 and
Theorem~10.1), the steps of Input~\ref{M-inp:density}, and the estimate with exponent $7/3$ on $[5/7, 16/21]$
(\lean{LinnikCGLDensity.cglDensity}); their Theorem~1.2 as they state it and their other results are not formalised.
The library's proof of their Theorem~10.1, a large-values bound with characters that they credit to Heath-Brown
\cite{hb79b}, follows the arguments of Kerr \cite{kerr} and Jutila \cite{jutila}; some lemmas of their paper were
formalised from an outline of their proofs; and the formalisation found six places where the text of their paper needs
a small repair, none affecting its results (\file{formalisations/chen-gupta-li-density.md}).

\subsection*{The statements}

The statement of Input~\ref{M-inp:density} and
those of the density estimates count the characters $\chi$ mod $q$, as elements of Mathlib's \lean{DirichletCharacter ℂ q}, for
which \lean{LFunction χ s = 0} at some $s$ in the rectangle. Mathlib's $L$-function of a character mod $q$ is the
analytic continuation of $\sum_n \chi(n) n^{-s}$, with the character taken mod $q$ and not reduced to a primitive
one. So $L(s, \chi)$ is the $L$-function of the primitive character inducing $\chi$ times the Euler factors at the
primes dividing $q$, whose zeros lie on $\operatorname{Re} s = 0$ and do not affect the counts for $\sigma \ge 1/2$.
In Lean the slack of Ingham's steps is $4\varepsilon$ for the $\varepsilon$ of
\lean{LinnikInghamC.charDensity\_ingham\_c}, whose exponent is $2 + 4c + 4\varepsilon$, with $c$ chosen so that
$2 + 4c = 3/(2 - \sigma_+)$, Ingham's exponent at the step's right end $\sigma_+$: $\varepsilon = 1/8000$ on the
first nineteen steps and $\varepsilon = 1/400000$ on the last. \file{formalisations/} records, for each result from the
literature, the differences between the Lean statement and the published one.
\subsection*{Independent checking of the inputs}

Independent checking of Input~\ref{M-inp:half-plane} is partial, and little of it is independent of OpenAI. OpenAI's
release has its own Lean formalisation of the main preprint's result. Its catalogue \texttt{lean/formalization.yaml}
\cite{openai-lean}, read on 25 October 2026, lists that preprint among the papers with a formalised main result, and
names the theorem \lean{OAI.DirichletCharacter.\allowbreak LFunction\_ne\_zero\_of\_seven\_eighths\_lt\_re}. Its
statement is for every modulus $q \ge 1$, every character mod $q$ and every $s$ with $\operatorname{Re} s > 7/8$,
against Mathlib's $L$-function, with the same exception at $s = 1$ for the trivial character, so it implies
Input~\ref{M-inp:half-plane} as stated here. The same catalogue gives the scope as partial progress and the review status
as unchecked, so it is OpenAI's own record and not an independent one. The author read that release's statements,
scanned its Lean and rebuilt it on 7 October 2026, which showed its two main theorems to depend on Lean's standard
three axioms only, but did not read its proofs in full. This library's proof of Input~\ref{M-inp:half-plane} is a new
formalisation of the companion preprint's route, on a different Lean and Mathlib; it imports nothing from the release
and was not compared with it. It is a machine-checked proof of the companion preprint's theorem, but the route is
OpenAI's and the formalisation is the author's, written with the AI assistance described in the paper's acknowledgements, so
it is not a review of the mathematics that is independent of OpenAI and of the author. The check recorded by the tracker AIProofWatch \cite{aiproofwatch} on 25 October 2026 concerns the
specialisation to the Riemann zeta function at $7/8$, not Dirichlet $L$-functions in general. There are third-party
reports of replaying OpenAI's Lean proof of the $7/8$ result, in at least one of them for the Dirichlet statement; they
check that proof, not the mathematics, and the author has not rerun them. So the author knows of no
review of the mathematics for Dirichlet $L$-functions in general that is independent of OpenAI and of the author.

For Input~\ref{M-inp:density} the author knows of no independent check either. The library's proof of Chen, Gupta and
Li's Proposition~3.1 is the author's formalisation of their text, written in the same way; some of their lemmas were
formalised from an outline of their proofs, and the formalisation found six places where the text needs a small repair
(\file{formalisations/chen-gupta-li-density.md}), which the author has sent to them. Lean's kernel checks that proof;
it does not check that its statement is theirs.

\subsection*{What the kernel checks}

The library's own sources, the directories \texttt{lib/core}, \texttt{lib/pnt}, \texttt{lib/dh},
\texttt{lib/zeta23} and \texttt{lib/check}, contain no \lean{sorry}, no declared axiom and no \lean{native\_decide},
which \texttt{scripts/check\_hygiene.py} checks. The library vendors 19 files of PrimeNumberTheoremAnd \cite{pnt-plus}
and, in \texttt{lib/vendor/zeta23}, 59 files of the \texttt{zeta23} folder of Anthropic's \texttt{formal-math}
\cite{formal-math}, each ported to the
library's toolchain by a patch that the repository records. That scan does not cover them; but the import cone of
Theorem~\ref{M-thm:main} contains 15 of the PrimeNumberTheoremAnd files and none of the \texttt{zeta23} ones, and not
\texttt{Wiener.lean}, the vendored file that contains \lean{sorry}. In any case the axiom check of the theorems in the
table of Appendix~\ref{M-app:claims} covers the whole dependency cone of each, and a \lean{sorry} in it would appear there as \lean{sorryAx}.

The certificate of Proposition~\ref{M-prop:certificate} is checked by kernel evaluation, \lean{decide +kernel}, and not by
compiled code, split into 5066 parts in 387 files: ranges of consecutive rows of the grid, and the heaviest $188$ rows in
pieces of columns. Each part proves that its sum is at most a bound $\lceil 1.001 \times \text{part} \rceil$, computed
exactly, and the bounds sum to $1151230848467235036 < 2^{60}$. Evaluating the parts tests about $0.48$ million,
$7.4$ million, $11.0$ million and $1.9$ million boxes of depths three to six below the decomposed cells, as
\file{tools/linnik/mirror2.py} counts them, besides the $27020$ candidate cells and $0.35$ million boxes for the kept
parts. The definitions the kernel evaluates are
arithmetic on Lean's core natural numbers, lists and booleans; their files import from Mathlib only the notation $\mathbb N$.

\subsection*{The artefact}

The library, the paper and this supplement are in the repository \cite{repo}; this version of the paper and of the
supplement describes the library at the commit that contains it. The toolchain is Lean \texttt{v4.35.0-rc3}, with Mathlib at commit \texttt{0f64d30a},
PrimeNumberTheoremAnd at \texttt{650d312} and \texttt{formal-math} at \texttt{fbdc36b}. The repository's README gives the
commands: \texttt{lake exe cache get} for Mathlib's prebuilt files, \texttt{scripts/build.sh} for the build, which
limits how many of the memory-hungry kernel checks of the certificate run at once, and
\texttt{python3 scripts/check\_axioms.py} for the axioms; the repository's continuous integration runs them on every
change to the library.
\section{Computations around the exponent}\label{sup:computations}

This section gives the computations behind the last two paragraphs of Section~\ref{M-sec:comparison} of the paper.
They are computations with the definitions of the certificate, by the scripts of \file{tools/linnik/}, and not theorems
of the library; nothing in the paper or the library depends on them.

\subsection*{The classical estimates in place of Chen, Gupta and Li's}

The gain does not depend on Chen, Gupta and Li's estimate. Take $g_7$ with its steps on $[5/7, 131/168]$ replaced by
steps of length $1/1680$ with the classical exponents and the slack $1/2000$ of the library's steps: Ingham's
$3/(2 - \sigma_+) + 1/2000$ up to $3/4$ and Huxley's $3/(3\sigma_- - 1) + 1/2000$ from $3/4$. The window is then set
at $3/4$, where both exponents are $12/5$, and the closed forms of Section~\ref{M-sec:ends} become $u_1 \ge 16/(5L) - 1$ and
$u_2 \le 1 - 6/(5L)$, which meet at $L = 11/5$. The search of \file{tools/linnik/density\_payoff.py} described below
(its variant \texttt{classical}) finds the window $[41098, 47082]/100000$ admissible at $L = 4537/2000 = 2.2685$, in
exact rationals, and the certificate's total on it $0.99835\mathcal K$ on $340$ cells a side; at $2.268$ it is
$1.00648\mathcal K$. So the sieve brings Huxley's $12/5$ down to about $2.27$, as it brings $7/3$ down to $2.224$,
and with Input~\ref{M-inp:half-plane} as the only input that is not refereed. This is a computation with the
certificate's definitions, not a theorem of the library: the library has no steps of Ingham's or Huxley's estimate of
length $1/1680$ on $[5/7, 131/168]$ (its Ingham steps end at $5/7$, and its Huxley steps, of length $1/280$, start at
$16/21$), and the certificate on this grid has not been checked by the kernel.
\subsection*{Below $278/125$}

The method as it stands does not go much further. At $278/125$ the certificate's total is $0.99754\mathcal K$
(Section~\ref{M-sec:certificate}); the kernel checks that it is below $0.99854\mathcal K$. Of it, about $0.971\mathcal K$
comes from the candidate cells that the sharp rule does not allow to be decomposed: about $0.507\mathcal K$ from those with
$\alpha_1 + \alpha_2$ beyond $1 - U_1/D$, whose cofactor is below the window, and about $0.464\mathcal K$, after the kept
parts, from those blocked only by their cells of depth three; the decomposed cells cost about $0.026\mathcal K$, so more
depth would not help. No refinement of the decomposition recovers most of the first part with this window: an
integer $p_1 p_2 m$ that those cells count with $\alpha_1 + \alpha_2 > 1 - U_1/D$ and $\alpha_1$ and $\alpha_2$ outside
$\mathcal U$ has no product of some of its prime factors with its exponent in $\mathcal U$, whatever the factorisation of
$m$ (a product with $p_1$ and another prime has exponent above $1 - U_1/D$, since the prime factors of $m$ are at least
$p_2$; one without $p_1$ that omits $p_2$ or a prime of $m$ has exponent at most $1 - \alpha_1 - \alpha_2 < U_1/D$; and
the others are $p_1$ and $p_2 m$, of exponents $\alpha_1$ and $1 - \alpha_1$), and $p_1 p_2$ is beyond the Type~I level $X'$. Such integers need a wider window, that is a better density exponent at $\sigma = 5/7$, or arithmetic
information of another kind. These figures, like the others of this subsection, are computations with the definitions of the
certificate (\file{tools/linnik/mirror2.py}), not theorems. A finer grid helps little: at $278/125$ the total is
$0.99167\mathcal K$ at $380$ cells a side and $0.98972\mathcal K$ at $400$ (and $0.99878\mathcal K$ at $360$, where the
cells lie less well against the window), but at $400$ cells a side $4447/2000 = 2.2235$ still fails, at
$1.00115\mathcal K$. The bounds of the cells are close to the integrals they bound: a constant $\omega(3) = 0.56438$
above $3$ in place of $0.565073$ would lower the total by only $6 \cdot 10^{-4}\mathcal K$.
The half-plane $\operatorname{Re} s > 7/8$ of OpenAI's main preprint \cite{openai-main}, in place of
Input~\ref{M-inp:half-plane}, would not help: the window condition (Definition~\ref{M-def:window}) would be needed only for
$\beta \le 7/8$, and the exponent of the kernel bound would become $b(\mu) = 3/8 + \mu/8 + \delta'(1 - \mu) + 2\varepsilon'$,
which is no larger for $\mu \le 1$; neither change reaches the corners that set the window's ends, on the first step at
$5/7$. Going lower would need more of the dropped terms evaluated or a better density exponent near $\sigma = 5/7$.

On the last, Chen, Gupta and Li note that the worst case for their estimate is $T = 1$ with no divisor $q_1$ of $q$ of a
suitable size, as for prime $q$, that the critical situation is then at $\sigma = 5/7$, and that improving the term which
dominates there would need a joint cancellation which ``seems quite difficult'' \cite[Remark~12.3]{cgl}. The script
\file{tools/linnik/density_payoff.py} measures what lower curves would give the method: for a curve made from $g_7$ it
finds a window admissible at $L$, checked at the corners in exact rationals, evaluates the certificate's total on it on
$340$ cells a side, and lowers $L$ in steps of $10^{-3}$ while the total stays below $0.999\mathcal K$, then tries
one further step of $5 \cdot 10^{-4}$ (\file{tools/linnik/README.md} describes the search). For $\mu \ge 1$ the window condition on a piece $A(1 - \beta)$ requires $L - \mu > A/2$, where $g_7 = 1$
counts as $A = 1/(1 - \beta)$, which tends to $7/3$ at $4/7$; so with the trivial count kept below $4/7$ and every
exponent capped at $30/13$ the script reaches only $2.2235$. The other curves replace the trivial count on $[11/20, 4/7]$
by one step of Ingham's estimate, of exponent $3/(2 - 4/7) + 10^{-5}$, and set the slack of the steps on $[4/7, 5/7]$ to
$10^{-5}$, both as \lean{LinnikInghamC.charDensity\_ingham\_c} allows (that change alone leaves $278/125$); with every
exponent then capped at $A$ the script reaches:
\begin{center}
\begin{tabular}{lccccc}
\toprule
cap $A$ & $2.33$ & $2.325$ & $2.32$ & $2.315$ & $30/13 \approx 2.3077$ \\
lowest $L$ & $2.2235$ & $2.2220$ & $2.2205$ & $2.2190$ & $2.2170$ \\
\bottomrule
\end{tabular}
\end{center}
These figures are partly those of $g_7$'s steps: for the capped curves the window's lower end is set not at $5/7$ but on
one of the steps of length $1/280$, whose exponent is taken at its left end (the first, $[607/840, 61/84]$, for the
cap $2.33$, and the last, $[163/210, 131/168]$, for the cap $30/13$).
With the same change and the exponent of each step inside $[7/10, 4/5]$ capped at Chen, Gupta and Li's bound at
$q_1 = q$ without its first term \cite[Theorem~1.2]{cgl}, the term that gives $4/(1 + \sigma)$, taken at the step's
left end plus $10^{-5}$, the script reaches $2.2090$; on $[7/10, 4/5]$ the largest of the exponents of the other three
terms is $15/(3 + 5\sigma)$, the exponent of Guth and Maynard's estimate for $\zeta$ \cite[Theorem~1.2]{guth-maynard}.
Capping the exponent on $[15/19, 11/12]$ at $2 + 10^{-5}$, as Heath-Brown's estimate
$\sum_{\chi \bmod q} N(\sigma, T, \chi) \ll (qT)^{2 - 2\sigma + \varepsilon}$ for $15/19 \le \sigma \le 1$ \cite{hb79}
would allow (the range as in the review of \cite{hb79} in zbMATH, Zbl~0393.10043), changes none of the figures for the four curves
on which the script tries it (the curve $g_7$ itself, the change on $[11/20, 5/7]$ alone, and the cap at $30/13$ and
the bound without the first term, both with that change). Nothing in the paper
or the library depends on these computations.


\begin{thebibliography}{99}
\raggedright
\bibitem{openai-main} OpenAI, \emph{The Quasi-Riemann Hypothesis: A Zero-Free Half-Plane $\operatorname{Re}(s) > 7/8$},
  preprint, 30 September 2026, \url{https://github.com/openai/math}, entry 003, read at commit \texttt{adc7f12}. Not
  refereed.
\bibitem{openai-lean} OpenAI, \emph{Catalog of papers with a formalized main result}, file
  \texttt{lean/formalization.yaml} (version v0.4) of the repository \texttt{openai/math}, a web page,
  \url{https://github.com/openai/math/blob/main/lean/formalization.yaml}, read 25 October 2026; with
  \texttt{lean/OAI/NumberTheory/DirichletL/Nonvanishing.lean} of the same repository. It is OpenAI's own record; it
  marks its scope as partial progress and its review status as unchecked.
\bibitem{aiproofwatch} AIProofWatch, record of OpenAI's \emph{The Quasi-Riemann Hypothesis: A Zero-Free Half-Plane
  $\operatorname{Re}(s) > 7/8$}, \url{https://aiproofwatch.com/updates/openai-math-003} (assessment reviewed 25 October 2026;
  reports an external Lean check of the Riemann zeta specialisation only).
\bibitem{cgl} B. Chen, V. Gupta and Y. C. Li, \emph{Large Value Estimates for Dirichlet Polynomials with Characters
  and Zero Density of Dirichlet $L$-Functions}, arXiv:2507.08296v2 (preprint; v2 of 27 July 2026, v1 of 11 July 2025
  was by B. Chen alone).
\bibitem{guth-maynard} L. Guth and J. Maynard, \emph{New large value estimates for Dirichlet polynomials},
  Ann. of Math. (2) 203 (2026), 623--675; arXiv:2405.20552.
\bibitem{hb79b} D. R. Heath-Brown, \emph{A large values estimate for Dirichlet polynomials}, J. London Math. Soc. (2) 20
  (1979), 8--18.
\bibitem{kerr} B. Kerr, \emph{Large values of Dirichlet polynomials and zero density estimates for the Riemann zeta
  function}, arXiv:1909.12075.
\bibitem{jutila} M. Jutila, \emph{Zero-density estimates for $L$-functions}, Acta Arith. 32 (1977), 55--62.
\bibitem{hb79} D. R. Heath-Brown, \emph{Zero density estimates for the Riemann zeta-function and Dirichlet
  $L$-functions}, J. London Math. Soc. (2) 19 (1979), 221--232.
\bibitem{brun} V. Brun, \emph{Le crible d'Eratosth\`ene et le th\'eor\`eme de Goldbach}, Videnskapsselskapets
  Skrifter I, Mat.-naturv. Klasse, Kristiania, 1920, no.~3.
\bibitem{pnt-plus} A. Kontorovich, T. Tao and contributors, \emph{PrimeNumberTheoremAnd}, a Lean project,
  \url{https://github.com/AlexKontorovich/PrimeNumberTheoremAnd}.
\bibitem{formal-math} Anthropic, \emph{formal-math}, a Lean repository, \url{https://github.com/anthropics/formal-math}.
\bibitem{repo} R. Comley, \emph{number-theory}, the repository of the paper, this supplement and its Lean library,
  \url{https://github.com/dickie81/number-theory}.
\end{thebibliography}

\end{document}
